\section*{Chapter \ref{ch:iv:probability-i} --- Probability I}

\begin{solution}{iq:iv:probability-i:1}
$5! = 120$. For AABBC, $5!/(2!\,2!) = 30$.
\iqlookfor{the product rule and division by the repeats.}
\end{solution}

\begin{solution}{iq:iv:probability-i:2}
$1 - (5/6)^4 = 671/1\,296 \approx 0.518$: slightly favourable at even money, an edge of about 3.5\% of the stake.
\iqlookfor{the complement and the translation into a bet's edge.}
\end{solution}

\begin{solution}{iq:iv:probability-i:3}
$\tfrac4{52} \times \tfrac3{51} = \tfrac1{221}$, or $\binom42/\binom{52}2 = 6/1\,326$.
\iqlookfor{sequential or combinatorial counting, with the same answer.}
\end{solution}

\begin{solution}{iq:iv:probability-i:4}
Stars and bars: $\binom{13}3 = 286$. With at least one each, give one to every venue and distribute the remaining
six: $\binom93 = 84$.
\iqlookfor{the multiset count and the shift for a minimum.}
\end{solution}

\begin{solution}{iq:iv:probability-i:5}
By exchangeability: A beats B with probability $\tfrac12$; A is best with $\tfrac16$; A beats both B and C when A is
first among the three, $\tfrac13$. Exhaustive enumeration of the 720 orders confirms all three.
\iqlookfor{symmetry, stated as an assumption about exchangeability.}
\end{solution}

\begin{solution}{iq:iv:probability-i:6}
$1 - \prod_{k<1000}(1 - k/10^6) \approx 1 - e^{-1000 \times 999/(2 \times 10^6)} \approx 0.393$. The chance reaches
one half when $n(n-1)/(2N) \approx \ln 2$, $n \approx \sqrt{2N\ln 2} \approx 1\,177$. Identifiers should be
sequential or wide (64 bits), not small random numbers.
\iqlookfor{the birthday approximation, the square-root scale, and the engineering conclusion.}
\end{solution}

\begin{solution}{iq:iv:probability-i:7}
$\frac{0.95 \times 0.002}{0.95 \times 0.002 + 0.02 \times 0.998} = \frac{95}{1\,093} \approx 0.087$. A good alert on a
rare event is mostly false alarms: in a population of 100\,000, 190 true alerts against 1\,996 false ones.
\iqlookfor{base rates, done by a table of counts or by odds.}
\end{solution}

\begin{solution}{iq:iv:probability-i:8}
If the report means ``the outcome contains at least one six'', condition on the 11 such outcomes: $\tfrac1{11}$.
If you looked at one die at random and it showed a six, the other die is independent: $\tfrac16$
(\cref{prop:iv:probability-i:protocols}). Ask: ``How did you come to tell me that: did you check both dice, or
look at one?'' A simulation of both protocols gives 0.091 and 0.167.
\iqlookfor{recognising that the answer depends on the protocol, and asking for it.}
\end{solution}

\begin{solution}{iq:iv:probability-i:9}
$53/5 = 10.6$. The four aces divide the other 48 cards into five gaps that are exchangeable, so each gap holds on
average $48/5$ cards, and the first ace follows the first gap. The symmetry replaces the sum over positions.
\iqlookfor{the gap \omterm{def:iv:probability-i:symmetry}{symmetry argument}.}
\end{solution}

\begin{solution}{iq:iv:probability-i:10}
In odds: prior $1:499$, each alert multiplies by $0.95/0.02 = 47.5$, so the posterior odds are $47.5^2/499 \approx
4.52$ and the probability $4.52/5.52 \approx 0.819$. Independence of the two alerts given the account's status is
the assumption that matters; two alerts built on the same data are not independent evidence.
\iqlookfor{odds updating, and the conditional-independence caveat.}
\end{solution}

\begin{solution}{iq:iv:probability-i:11}
Weight each envelope by the chance of drawing a winning ticket first: $\tfrac14 \times (1, \tfrac12, 0,
\tfrac23)$. Given a winning first draw, the chance that the second is winning is 1 from the first envelope, 0 from
the second, and $\tfrac12$ from the fourth (one winning and one losing remain). The answer is $(1 \times 1 +
\tfrac12 \times 0 + \tfrac23 \times \tfrac12)/(1 + \tfrac12 + \tfrac23) = \tfrac43 / \tfrac{13}6 = \tfrac8{13} \approx
0.615$. This is a variant of Bertrand's box; the trap is to answer $\tfrac12$ by counting envelopes rather than
tickets.
\iqlookfor{weighting hypotheses by the likelihood of the observation.}
\end{solution}

\begin{solution}{iq:iv:probability-i:12}
Staying wins with probability $\tfrac14$. Switching: with probability $\tfrac34$ the prize is behind one of the
other three doors, and after the host opens an empty one it is behind one of the two you may switch to; choosing
at random between them wins half the time: $\tfrac34 \times \tfrac12 = \tfrac38$. Switching is better. An
enumeration over the prize's position, the host's choice and yours confirms it.
\iqlookfor{conditioning on the host's informed choice, and the general principle behind the classic three-door
answer.}
\end{solution}

\begin{solution}{iq:iv:probability-i:13}
In the 60-by-60 square of arrival times, the orders are more than ten seconds apart in two triangles of legs 50:
area $50^2 = 2\,500$ of $3\,600$. The chance of arriving within ten seconds is $1 - (50/60)^2 = 11/36 \approx
0.306$.
\iqlookfor{a geometric picture of a joint uniform law.}
\end{solution}

\begin{solution}{iq:iv:probability-i:14}
Each position has 14 equally likely states (profitable or not, times seven days). The outcomes in which at least
one position is profitable and opened on Monday number $14^2 - 13^2 = 27$; those in which both are profitable
and at least one of them was opened on Monday number $7^2 - 6^2 = 13$. The answer is $13/27$. It is not
$\tfrac13$ (the answer to ``at least one is profitable'') because the Monday detail makes the case of two profitable
positions nearly twice as likely to produce the report: the more specific the information, the closer the answer
to $\tfrac12$.
\iqlookfor{careful counting of the conditioning event, and the reason extra detail changes it.}
\end{solution}
